%==============================================================================
%  A Visualisation Layer for Web-Based Digital Twin Platforms
%  Techniques, Architecture and an Implementation Plan for DTaaS
%
%  Build:  pdflatex main && bibtex main && pdflatex main && pdflatex main
%  (or: latexmk -pdf main.tex)
%==============================================================================
\documentclass[11pt,a4paper]{article}

%------------------------------------------------------------------------------
% Packages
%------------------------------------------------------------------------------
\usepackage[utf8]{inputenc}
\usepackage[T1]{fontenc}
\usepackage{lmodern}
\usepackage[english]{babel}

\usepackage{amsmath}
\usepackage{amssymb}
\usepackage{textcomp}

\usepackage[a4paper,margin=2.4cm,headheight=14pt]{geometry}
\usepackage{microtype}
\usepackage{booktabs}
\usepackage{longtable}
\usepackage{tabularx}
\usepackage{array}
\usepackage{multirow}
\usepackage{graphicx}
\usepackage[table]{xcolor}
\usepackage{enumitem}
\usepackage{listings}
\usepackage{fancyhdr}
\usepackage{titlesec}
\usepackage{caption}
\usepackage{natbib}
\usepackage[hidelinks]{hyperref}

%------------------------------------------------------------------------------
% Colours
%------------------------------------------------------------------------------
\definecolor{dtblue}{HTML}{0F6C8A}
\definecolor{dtamber}{HTML}{A4670F}
\definecolor{dtrule}{HTML}{B4C3CB}
\definecolor{dtgrey}{HTML}{4D606C}
\definecolor{codebg}{HTML}{F2F5F7}

\hypersetup{
  colorlinks=true,
  linkcolor=dtblue,
  citecolor=dtblue,
  urlcolor=dtblue,
  pdftitle={A Visualisation Layer for Web-Based Digital Twin Platforms},
  pdfsubject={Techniques, architecture and implementation plan for the DTaaS client}
}

%------------------------------------------------------------------------------
% Headers / footers
%------------------------------------------------------------------------------
\pagestyle{fancy}
\fancyhf{}
\fancyhead[L]{\small\textsc{DTaaS Visualisation Layer}}
\fancyhead[R]{\small\thepage}
\renewcommand{\headrulewidth}{0.4pt}

%------------------------------------------------------------------------------
% Section styling
%------------------------------------------------------------------------------
\titleformat{\section}{\normalfont\Large\bfseries\color{dtblue}}{\thesection}{0.7em}{}
\titleformat{\subsection}{\normalfont\large\bfseries}{\thesubsection}{0.6em}{}
\titleformat{\subsubsection}{\normalfont\normalsize\bfseries}{\thesubsubsection}{0.6em}{}

%------------------------------------------------------------------------------
% Listings
%------------------------------------------------------------------------------
\lstdefinestyle{ts}{
  basicstyle=\ttfamily\footnotesize,
  backgroundcolor=\color{codebg},
  keywordstyle=\color{dtblue}\bfseries,
  commentstyle=\color{dtgrey}\itshape,
  stringstyle=\color{dtamber},
  showstringspaces=false,
  breaklines=true,
  frame=leftline,
  rulecolor=\color{dtrule},
  framesep=6pt,
  xleftmargin=8pt,
  numbers=none,
  tabsize=2,
  morekeywords={interface,type,export,import,const,readonly,string,number,boolean,void,Promise,null,extends,implements,function,return,from,as}
}
\lstset{style=ts}

%------------------------------------------------------------------------------
% Table helpers
%------------------------------------------------------------------------------
\newcolumntype{L}[1]{>{\raggedright\arraybackslash}p{#1}}
\newcolumntype{C}[1]{>{\centering\arraybackslash}p{#1}}

\newcommand{\pkg}[1]{\texttt{#1}}
% Breakable typewriter text, for long file names in narrow table columns
\newcommand{\fn}[1]{{\footnotesize\path{#1}}}
\newcommand{\tech}[1]{\textbf{T#1}}
\newcommand{\req}[1]{\textbf{R#1}}

% Callout box
\newenvironment{keypoint}
  {\par\vspace{0.5em}\noindent\begin{tabular}{@{}p{0.3em}@{\hspace{0.8em}}p{0.92\linewidth}@{}}
   \cellcolor{dtamber} & \itshape}
  {\end{tabular}\par\vspace{0.5em}}

\setlist[itemize]{itemsep=2pt, topsep=4pt}
\setlist[enumerate]{itemsep=2pt, topsep=4pt}

\captionsetup{font=small, labelfont={bf,color=dtblue}}


\title{\bfseries A Visualisation Layer for Web-Based\\ Digital Twin Platforms\\[0.3em]
\large Problem, Solution Architecture, Techniques, Open-Source Packages\\
and an Implementation Plan for the DTaaS Client}

\author{Prepared for the INTO-CPS Association \\ Digital Twin as a Service (DTaaS) project}

\date{13 September 2026}

\begin{document}

\maketitle
\thispagestyle{fancy}

\begin{abstract}
\noindent
The DTaaS web client has no spatial visualisation capability: it is a React~19 / Vite
application whose dependency set contains no 3D, image-overlay or
scientific-visualisation library. Any digital twin (DT) with a spatial subject can be
executed, stored and logged on the platform, but seen only as a time series. This report
states the problem (\S\ref{sec:problem}), proposes a layered architecture in which the
\emph{binding layer} --- not the renderer --- is the reusable asset
(\S\ref{sec:architecture}), catalogues fourteen visualisation techniques mapped onto the
domains DTaaS serves (\S\ref{sec:techniques}), names an open-source package for each with
version and licence (\S\ref{sec:packages}), and sets out a phased plan for a React
component package, \pkg{@dtaas/visualisation} (\S\ref{sec:implementation}).

Three claims organise it. First, such a package must not be a renderer: renderers are
domain-specific and short-lived, whereas anchor resolution, the temporal state store and
the encoding vocabulary are durable. Second, a 3D scene is one \emph{substrate} among
several --- photographs and live video are equally legitimate places to put simulation
output, and are what many users ask for first. Third, roughly a third of what such a
layer normally builds is already deployed: \pkg{dtaas-services} provides RabbitMQ (AMQP
and MQTT from one broker), InfluxDB, ThingsBoard, Grafana, MongoDB and GitLab, and the
client already ships \pkg{@gitbeaker/rest} and \pkg{react-iframe}. The plan consumes
those rather than duplicating them, and makes visualisation an additional DTaaS library
asset type so a configuration can be authored, versioned in Git, reviewed as a merge
request, shared and reused exactly as any other DT asset is.
\end{abstract}

\section{Introduction}
\label{sec:intro}

Digital Twin as a Service (DTaaS) is a platform for the creation, execution and sharing
of digital twins, owned by the INTO-CPS Association~\citep{talasila2024composable}. Its
monorepo comprises a React web client, library, runner and logger microservices, and the
\pkg{dtaas-services} service manager. Its defining idea is \emph{composability}: a twin
is assembled from reusable library assets --- data, models, functions and tools ---
selected in the client and executed by the runner.

Three existing capabilities change what a visualisation layer must build rather than
simply use. \textbf{A full data stack}: \pkg{dtaas-services} installs, in TLS mode,
InfluxDB, Grafana, RabbitMQ as an AMQP broker \emph{with its MQTT plugin enabled} --- so
one broker serves both protocols --- MongoDB, ThingsBoard (IoT device management, with a
PostgreSQL backend) and GitLab as OAuth2 identity provider and Git repository service.
\textbf{Git in the browser}: the client already depends on \pkg{@gitbeaker/rest}~43.8, so
it can read \emph{and write} repository files under the same OAuth2 session, meaning
versioned persistence needs no new backend. \textbf{Real Linux workspaces}: each user
gets a container exposing Jupyter, VS~Code, \pkg{ungit} and an xfce desktop over VNC
\emph{in the browser}, embedded through \pkg{react-iframe}; files a twin execution
produces are addressable by URL.

Against that, the client's own dependencies are React~19.2, Vite~8, TypeScript~6, MUI~9,
Redux Toolkit~2, React Router~7 and Monaco, with routes for \texttt{account},
\texttt{auth}, \texttt{config}, \texttt{digitaltwins}, \texttt{library},
\texttt{measurement} and \texttt{workbench} --- and \emph{no} 3D, image-overlay or
scientific-visualisation dependency at all.

\paragraph{Goal.} A reusable React component package, \pkg{@dtaas/visualisation}, that
(a) accepts visual substrates from arbitrary URLs --- 3D assets, photographs, floor plans
and live video --- including assets served by the library microservice, GitLab or a
workspace; (b) drives their visual state from heterogeneous live streams, reusing the
brokers and databases already deployed; (c) serves several domains (city, machines,
buildings, hospitals, and the marine and mobility domains evidenced in the literature
reviewed) without a rebuild for each; and (d) couples \emph{simulated} to
\emph{measured} state in a single view --- the capability that distinguishes a digital
twin from a monitoring dashboard.

\paragraph{Two requirements that are easy to under-weight.} A large share of what users
ask for is not a 3D model. They have a \emph{photograph} --- a traffic camera still, a
room, a plant floor, a scanned plan --- and want simulation and live readings drawn on
it, correctly placed (technique \tech{13}). Where a camera exists they want the
\emph{live feed beside the simulation}, moving together in time, so divergence is
directly observable (technique \tech{14}). Both are treated here as first-class
substrates, which shapes the architecture in Section~\ref{sec:architecture}.

\paragraph{Out of scope.} Dashboarding, which Grafana and ThingsBoard already do well ---
the right move is to \emph{drive} them from the same clock, not rebuild them. And asset
authoring: the package consumes substrates. The exceptions are binding authoring and
image registration, because both bind a signal to a location and no external tool can do
that for us.

\paragraph{Sources.} This report consolidates a survey of commercial platforms (Autodesk
Tandem and Platform Services, AWS IoT TwinMaker, Bentley iTwin.js, Samsara, NVIDIA
Omniverse/OpenUSD); an audit of open-source packages with versions and licences verified
against the npm registry and GitHub API on 12--13 September 2026; and eight peer-reviewed
papers supplied in \texttt{sources/}. Appendix~\ref{app:sources} records provenance,
including three supplied files that are failed downloads rather than papers.

\section{Problem}
\label{sec:problem}

\begin{quote}\itshape
A web-based DT platform must render arbitrary visual substrates --- 3D assets,
photographs and live video --- loaded from arbitrary URLs, driven by heterogeneous live
data streams, across unrelated application domains, without rebuilding the stack each
time a new domain or medium appears. No open-source package solves this: every candidate
solves one domain deeply or all domains shallowly, and almost none treat an image or a
video feed as a legitimate place to put simulation output.
\end{quote}

\subsection{Six difficulties}

\textbf{P1 Domain rebuild.} 3D libraries are domain-shaped --- a BIM viewer models
storeys and IFC property sets, a geospatial engine models tiles and terrain, a robotics
viewer models joints. Serving four domains means either one poor fit in three of them or
four integrations to maintain. \citet{rundel2023leveraging} report the same pressure at
city scale: operational city DTs with simulation (Virtual Singapore, Herrenberg) were
\emph{impracticable} for small municipalities --- Virtual Singapore alone exceeded
USD~73~million over five years --- and their answer was procedural generation from
existing public geodata rather than bespoke authoring. The lesson generalises: derive
the scene from data that already exists.

\textbf{P2 Anchoring.} Something must connect a signal to a location. TwinMaker binds
entity/component/property onto a GLB submodel; iTwin.js keys on Element~ID; Tandem
classifies assets explicitly. All three \emph{author} the link and persist it; runtime
inference from coordinates or names breaks on re-export. DTaaS has an advantage here:
ThingsBoard already holds a device and asset registry, so the anchor namespace need not
be invented.

\textbf{P3 Sustained update.} Applying an encoding once is trivial; applying it to
$10^4$--$10^6$ elements at 10--60~Hz without rebuilding the scene graph is not. This
constrains technique choice far more than aesthetics.

\textbf{P4 Heterogeneity.} \citet{ali2024spatial} name this as open research: combining
3D geometry, time series and map layers legibly, and visualising correlations across
space and time, ``needs more research''. Part of the problem is genuinely unsolved, not
merely unbuilt.

\textbf{P5 Simulation coupling.} The problem commercial platforms serve worst. Tandem,
TwinMaker and Samsara visualise \emph{measurement} at high polish; simulation lives in
another tool with another viewer. Yet the residual between predicted and measured state
is the high-value signal, because it represents a physics violation rather than a
statistical oddity. \citet{liao2025digital} co-visualise sensed and simulated
deformation of a soft robot in one AR scene, reporting errors of 0--7.14\%.

\textbf{P6 Substrate plurality.} Most libraries assume the answer is a 3D scene. In
practice a large share of requests are a \emph{photograph} (traffic still, room, plant
floor, scanned plan) on which simulation should be drawn; a \emph{live video feed} to be
shown beside the simulation; or a \emph{rendered output file} from a workspace execution
that needs only placing and time-aligning. Treating these as second-class leaves the
visualisation layer unusable until someone commissions a 3D model.

\begin{keypoint}
DTaaS is unusually placed on P5: its INTO-CPS lineage means FMI co-simulation already
produces the predicted channel every surveyed commercial platform lacks, and InfluxDB
already provides somewhere to put both channels. The cost of exploiting this is small
\emph{provided} measured and simulated state are two channels of one signal from the
start.
\end{keypoint}

\subsection{Why the obvious approaches fail}

\textbf{One 3D library} fails on P1 and P6. Pascal (\pkg{@pascal-app/*}, MIT) gives a
full architectural editor but its document model is sites\,$\rightarrow$\,buildings\,%
$\rightarrow$\,levels\,$\rightarrow$\,walls, in which a pump, a fjord or a photograph is
not a citizen, and its viewer exposes no live-data binding. \pkg{@thatopen/components}
has the best data semantics of any BIM option and is likewise BIM-shaped. CesiumJS is
decisive at city scale and clumsy for one machine. \textbf{Everything on three.js} fails
on effort, and is the wrong tool for annotating a JPEG. \textbf{A game engine} works ---
\citet{rundel2023leveraging} use Unity+SUMO via TraCI, \citet{vasilijevic2024trondheim}
Unity+Cesium --- but produces an artefact hosted and versioned outside the platform; its
proper home in DTaaS is the workspace desktop (\S\ref{sec:workspace}).
\textbf{Grafana and ThingsBoard alone} solve charting completely and the spatial problem
not at all. \textbf{A commercial platform} does not solve P5, and only iTwin.js (MIT) is
installable.

\subsection{Requirements}

\begin{longtable}{@{}L{1.1cm}L{13.4cm}@{}}
\toprule
\textbf{ID} & \textbf{Requirement} \\
\midrule
\endhead
\req{1} & \textbf{Substrate independence.} Adding a domain or a medium means adding an
adapter behind a stable interface, not forking the package. \\
\req{2} & \textbf{Authored anchoring.} Signal-to-location bindings are explicit,
persisted, versionable and human-editable. \\
\req{3} & \textbf{Transport agnosticism.} MQTT and AMQP (both via the installed
RabbitMQ), ThingsBoard telemetry, InfluxDB, SQL, HTTP and files normalise to one stream
shape. \\
\req{4} & \textbf{Temporal addressing.} Every encoding reads state at a playhead, never
``latest value''. \\
\req{5} & \textbf{Dual-channel state.} Measured and simulated values of one signal are
first-class and separately addressable. \\
\req{6} & \textbf{Sustained update.} $O(\text{changed elements})$ per frame, no
scene-graph rebuild. \\
\req{7} & \textbf{Composability as a DTaaS asset.} A visualisation is storable in the
library, selectable in composition, shareable between users. \\
\req{8} & \textbf{Licence cleanliness.} MIT, Apache-2.0 or BSD only; network copyleft is
disqualifying for a platform client. \\
\req{9} & \textbf{Graceful degradation.} A twin with no 3D asset still gets a usable
view; a browser without WebGPU still renders. \\
\req{10} & \textbf{Reuse of the installed stack.} Consume InfluxDB, RabbitMQ,
ThingsBoard, Grafana and GitLab as deployed; no parallel broker, historian, dashboard or
file store. \\
\req{11} & \textbf{Media synchronisation.} Where video and simulation are shown
together, their time alignment is explicit, measurable and correctable. \\
\bottomrule
\end{longtable}

\noindent
\req{8} eliminates \pkg{xeokit-sdk} and \pkg{realvirtual WEB} (both AGPL-3.0) and flags
\pkg{Potree} and \pkg{loaders.gl} (\texttt{NOASSERTION}) for checking. \req{10} is what
most changes the plan relative to a greenfield design.


\section{Solution Architecture}
\label{sec:architecture}

\subsection{The invariant beneath every surveyed platform}

Strip the branding from Tandem, TwinMaker, iTwin.js and Omniverse and the same three
layers appear. \textbf{Anchor}: a stable identifier tying a signal to a location on a
substrate --- IFC GUID, glTF node, USD prim, Element~ID, or a point or polygon in image
coordinates. \textbf{Encode}: a rule turning a value into a visual property --- colour,
opacity, visibility, transform, flow speed, ghost offset, glyph, label; this is where
commercial platforms put their product, and TwinMaker's no-code rule editor \emph{is}
what customers pay for. \textbf{Sustain}: applying encodings continuously at scale
without rebuilding the scene; iTwin.js devotes a native tile generator and a WebGL
frontend to this layer alone.

\begin{keypoint}
The renderer is not in this list. It is an implementation detail of the \emph{sustain}
layer. Treating it as the centre of the architecture is what makes a visualisation
package domain-locked --- and what makes packages unable to treat a photograph or a
video frame as a place where data can live.
\end{keypoint}

\subsection{Six layers}

Layers~1--4 are domain- and substrate-independent and form the reusable core; layer~5 is
a thin per-substrate adapter; layer~6 is third-party.

\paragraph{Layer 1 --- Transport.} Every source normalises to one event shape, with
optional historical backfill:

\begin{lstlisting}
export type Channel = 'measured' | 'simulated' | 'predicted' | 'setpoint';

export interface SignalSample {
  readonly twinId: string; readonly signalPath: string;   // "ward3/ahu1/supplyTemp"
  readonly channel: Channel; readonly ts: number;         // epoch ms, source-assigned
  readonly value: number | string | boolean | number[];
  readonly quality?: 'good' | 'stale' | 'bad';
}

export interface TransportAdapter {
  readonly id: string;
  connect(cfg: unknown): Promise<void>;
  subscribe(paths: string[], sink: (s: SignalSample) => void): () => void;
  range?(paths: string[], from: number, to: number): Promise<SignalSample[]>;
  disconnect(): Promise<void>;
}
\end{lstlisting}

Adapters follow the installed stack, not a wish-list: MQTT over WSS and AMQP over
Web-STOMP, both against the single RabbitMQ; ThingsBoard telemetry over its WebSocket
API; InfluxDB for \texttt{range()}; SQL/HTTP for the logger and arbitrary DT endpoints;
and a file adapter replaying CSV or Parquet from the library, GitLab or a workspace
URL --- which makes the stack testable offline and lets an FMI output file serve as the
\texttt{simulated} channel with no runtime integration.

\paragraph{Layer 2 --- Temporal state store, hot and cold.} A ring buffer per
\texttt{(signalPath, channel)} holds the live window; when the playhead moves outside
it, the store issues an InfluxDB \texttt{range()} query and serves from the returned
window. Every consumer calls \texttt{valueAt(path, channel, t)}, never ``current
value''. This implements \req{4} and unlocks 4D replay (\tech{6}), residuals (\tech{9}),
image-overlay replay (\tech{13}) and video sync (\tech{14}) at low marginal cost. The
playhead, subscriptions and connection status belong in Redux Toolkit; the sample
buffers deliberately do not, because dispatching an action per sample destroys frame
rate.

\begin{keypoint}
Having InfluxDB deployed converts the time store from a hard problem into a cache.
Without a historian, honest replay means holding everything in memory or building a
backend; with one, it is a windowed query behind an interface the rest of the package
cannot distinguish from live data.
\end{keypoint}

\paragraph{Layer 3 --- Anchor resolution.} Anchor kinds span three media, plus the
imported registry:

\begin{lstlisting}
export type AnchorKind =
  | 'ifc-guid' | 'gltf-node' | 'usd-prim' | 'instance-id'   // 3D substrates
  | 'geo'                                                    // georeferenced
  | 'image-point' | 'image-region' | 'image-plane'           // image substrates
  | 'video-panel'                                            // video substrates
  | 'tb-entity';                                             // ThingsBoard device/asset
\end{lstlisting}

An anchor carries \texttt{signalPath}, \texttt{kind}, \texttt{ref}, a human
\texttt{label} and, for image substrates, the \texttt{homographyId} of the calibration
that applies. \texttt{tb-entity} matters: because ThingsBoard holds the registry, an
anchor can pair \emph{this device} with \emph{that IFC element}, so the signal namespace
is imported rather than invented and a rename does not silently break the binding.

\paragraph{Layer 4 --- Encoding rules.} Declarative, serialisable, and the genuinely
reusable content of the package: \texttt{colorScale}, \texttt{visibility},
\texttt{transform}, \texttt{flow}, \texttt{residual}, \texttt{ghost}, \texttt{label} and
\texttt{attention} for geometric substrates; \texttt{glyph}, \texttt{regionFill},
\texttt{fieldOverlay} and \texttt{trajectory} for image and video substrates. Because
encodings are plain data, a complete visualisation is a JSON document --- which is what
makes \req{7} achievable and, given \pkg{@gitbeaker/rest}, what makes it storable and
reviewable with no new backend.

\paragraph{Layer 5 --- Substrate adapters.} The seam that delivers \req{1}. Generalising
from \emph{renderer} to \emph{substrate} is what admits images and video as peers of 3D:

\begin{lstlisting}
export interface SubstrateAdapter {
  readonly id: string; readonly supports: AnchorKind[];
  mount(container: HTMLElement, descriptor: SubstrateDescriptor): Promise<void>;
  resolve(anchor: Anchor): ElementRef | null;
  apply(ref: ElementRef, enc: ResolvedEncoding): void;   // hot path, must be O(1)
  frame(ref: ElementRef, opts?: FrameOptions): void;
  pick(x: number, y: number): ElementRef | null;
  setTime(t: number): void;                              // media seek, animation
  dispose(): void;
}
\end{lstlisting}

\begin{longtable}{@{}L{1.9cm}L{5.2cm}L{7.4cm}@{}}
\toprule
\textbf{Adapter} & \textbf{Backing library} & \textbf{Substrate and domains} \\
\midrule
\endhead
\texttt{image} & \pkg{react-konva}; \pkg{OpenCV.js} for calibration; \pkg{Leaflet}
(\texttt{CRS.Simple}) for very large images & Photographs, floor plans, scanned
drawings, rendered output frames. All domains; cheapest substrate, build first. \\
\addlinespace
\texttt{aec} & \pkg{@thatopen/components} + \pkg{@thatopen/fragments} & IFC models:
hospitals, buildings, plants. \\
\addlinespace
\texttt{mesh} & \pkg{react-three-fiber} + \pkg{drei} (+ \pkg{urdf-loader}) & glTF
assemblies: machines, robots, equipment. \\
\addlinespace
\texttt{video} & \texttt{HTMLVideoElement} + \pkg{hls.js} or WebRTC & Live and recorded
camera feeds beside or beneath simulation output. \\
\addlinespace
\texttt{geo} & \pkg{CesiumJS} / \pkg{3DTilesRendererJS} (+ \pkg{deck.gl}) & City,
mobility, marine; anything georeferenced or tiled. \\
\addlinespace
\texttt{field} & \pkg{vtk.js} & Simulation output: volumes, isosurfaces, streamlines.
Own canvas, synchronised camera. \\
\addlinespace
\texttt{embed} & \pkg{react-iframe} & Grafana panels, ThingsBoard dashboards and the
workspace VNC desktop, driven by the shared playhead. \\
\bottomrule
\end{longtable}

\subsection{Mapping onto the deployed DTaaS stack}
\label{sec:stack}

\begin{longtable}{@{}L{2.5cm}L{3.9cm}L{8.0cm}@{}}
\toprule
\textbf{Service} & \textbf{Role} & \textbf{How the package uses it} \\
\midrule
\endhead
\textbf{RabbitMQ} & Live transport, AMQP + MQTT from one broker & Two Layer-1 adapters
against one deployment: \pkg{mqtt.js} over WSS, \pkg{@stomp/stompjs} over Web-STOMP. No
second broker. \\
\addlinespace
\textbf{InfluxDB} & Historian & Cold tier of the Layer-2 store via
\pkg{@influxdata/influxdb-client}; makes \tech{6} and \tech{9} possible over arbitrary
past windows. \\
\addlinespace
\textbf{ThingsBoard} & Device/asset registry; telemetry & Anchor namespace
(\texttt{tb-entity}); a second live transport; dashboards embeddable. \\
\addlinespace
\textbf{Grafana} & Charting & \emph{Not} reimplemented. Embedded and driven by the
shared playhead: the viewer writes the current window into the panel URL's
\texttt{from}/\texttt{to}, so scrubbing the substrate moves the chart. \tech{10} for
near-zero cost. \\
\addlinespace
\textbf{GitLab} & Versioned storage and identity & Visualisation assets, anchor maps and
(via LFS) models and photographs read and written with \pkg{@gitbeaker/rest} under the
existing OAuth2 session. A threshold change becomes a merge request. \\
\addlinespace
\textbf{MongoDB} & Document store & Optional: large generated anchor maps, per-user view
state that should not pollute Git history. \\
\addlinespace
\textbf{Library / runner} & Files; DT execution & Resolves \texttt{lib://} substrates
(needs MIME types for \texttt{.glb}, \texttt{.ifc}, \texttt{.frag}); produces the
\texttt{simulated} channel and rendered frames. \\
\bottomrule
\end{longtable}

\begin{keypoint}
The highest-value integration in this table is the cheapest: driving an embedded Grafana
panel from the visualisation playhead. It gives linked time-and-space navigation without
writing a charting component, and respects the division of labour the marine DT pilot
arrived at independently~\citep{vasilijevic2024trondheim}.
\end{keypoint}

\subsection{Persistence, the workspace, and composability}
\label{sec:workspace}

Because \pkg{@gitbeaker/rest} is already a client dependency and GitLab is already the
identity provider, the visualisation asset needs no bespoke storage: the viewer reads
\texttt{visualisation.json} and its substrate at a given ref, the binding editor commits
changes back, review happens as a merge request diff, and large substrates live in Git
LFS beside the configuration referencing them.

Some visualisation will never happen in a browser canvas --- full CFD post-processing,
ParaView, Blender, \texttt{rviz}, a Unity city build. All already run in the DTaaS
workspace, whose xfce desktop is reachable over VNC inside the client. The architecture
should embrace this on three conditions: the desktop mounts through the \texttt{embed}
adapter like any other substrate; output files written to the workspace come back by URL
and are consumed by the \texttt{image} or \texttt{mesh} adapter, so a heavyweight render
becomes a navigable, time-aligned artefact for everyone else; and where a workspace tool
cannot accept the client's playhead, the layout says the pane is independent rather than
implying synchrony. This gives an honest answer to ``can DTaaS do proper scientific
visualisation?'' --- yes, in the workspace --- without pretending the web layer will
match ParaView.

Finally, DTaaS already treats data, models, functions and tools as composable library
assets. The consistent move is to add a \emph{visualisation} asset type, so a
\texttt{visualisation.json} sits beside the models and functions of a twin and is picked
in the same composition flow. An abbreviated example spanning three substrates:

\begin{lstlisting}
{ "schemaVersion": "1.1", "name": "junction-7-live",
  "layout": { "type": "split", "panes": ["photo", "camera", "grafana"] },
  "substrates": {
    "photo":   { "adapter": "image", "source": "git://twins/junction7/junction7.jpg",
                 "calibration": { "type": "homography", "id": "h1",
                   "imagePoints": [[412,880],[1502,872],[1710,1040],[210,1050]],
                   "worldPoints": [[0,0],[24,0],[24,12],[0,12]], "units": "m" } },
    "camera":  { "adapter": "video", "source": ".../junction7/index.m3u8",
                 "latencyHintMs": 2500 },
    "grafana": { "adapter": "embed", "source": "https://grafana.example/d/abc/junction7",
                 "timeParams": ["from", "to"] } },
  "transports": [
    { "adapter": "mqtt", "url": "wss://rabbitmq.example:15676/ws",
      "topics": ["dtaas/junction7/+/count"], "channel": "measured" },
    { "adapter": "influx", "org": "dtaas", "bucket": "junction7", "role": "history" },
    { "adapter": "file", "url": "lib://data/junction7-sumo-out.csv",
      "channel": "simulated" } ],
  "anchors": [
    { "signalPath": "junction7/armA/queueLen", "kind": "image-region",
      "ref": "poly:armA", "homographyId": "h1", "label": "North arm queue" } ],
  "encodings": [
    { "target": "junction7/armA/queueLen", "substrate": "photo",
      "encoding": { "type": "regionFill", "domain": [0,40], "scheme": "traffic" } },
    { "target": "junction7/armA/queueLen", "substrate": "photo",
      "encoding": { "type": "residual", "against": "simulated", "domain": [-8,8] } } ] }
\end{lstlisting}

A visualisation therefore becomes shareable (replace the photograph and the homography),
reviewable (JSON in Git) and executable by the runner (a headless adapter emits the same
overlay as an image file during a DT execution).

\subsection{Design rules}

\begin{enumerate}[itemsep=1pt]
  \item The binding layer is the product; substrates are dependencies. Layers~1--4 must
        import no rendering, imaging or media library --- enforce with a lint rule.
  \item One clock. Scene, image overlay, video element and embedded Grafana panel all
        sample the playhead.
  \item Encodings are data. Anything expressible only as code is a gap in the
        vocabulary.
  \item Hot path outside Redux. Redux holds control state; per-frame sampling reads
        buffers directly.
  \item Reuse before build. If InfluxDB, RabbitMQ, ThingsBoard, Grafana or GitLab does
        it, call it.
  \item Every adapter lazily loaded. An image twin must never ship the BIM engine.
\end{enumerate}


\section{Visualisation Techniques by Domain}
\label{sec:techniques}

Fourteen techniques, ordered by depth of coupling between physical asset, simulation and
live stream. \tech{1}--\tech{12} apply mainly to geometric substrates; \tech{13} and
\tech{14} cover image and video.

\begin{description}[leftmargin=1.55cm, style=nextline, itemsep=3pt, font=\bfseries\color{dtblue}]

\item[T1 Semantic anchoring] Bind \texttt{signalPath}\,$\rightarrow$\,element using
identifiers already in the model. Everything downstream is a lookup on this map; the
inverse lookup (picking) makes the scene interrogable. \emph{All domains ---
precondition, not option.}

\item[T2 Feature overrides] Recolour, fade or hide individual elements via a per-element
attribute read by one shared shader, keeping cost proportional to \emph{changed}
elements. iTwin.js \texttt{FeatureOverrideProvider}; TwinMaker model shader.
\textbf{Highest-leverage technique: \tech{3}, \tech{9} and \tech{12} all consume it.}

\item[T3 Sparse-sensor field mapping] Twenty sensors, one continuous surface.
Interpolate (IDW, RBF, nearest) into a texture sampled by the surface shader. The
interpolation is the honest part and must expose its uncertainty. \emph{Hospitals,
buildings, marine.}

\item[T4 Volumetric simulation fields] CFD, thermal, acoustic or dispersion results as
isosurfaces, cutting planes or ray-marched volumes registered to the sensor geometry.
Absent from Tandem, TwinMaker and Samsara; central to
\citet{vasilijevic2024trondheim}, who render OpenDrift and DREAM forecasts into the
twin. \emph{Marine, city, buildings, machines.}

\item[T5 Flow advection] Rate and direction as movement --- particles in a pipe, arrows
in a duct, vehicles on a link. Distinguishes ``flowing slowly'' from ``stopped'', which
a colour ramp cannot. TwinMaker ships it as the \emph{motion indicator}. \emph{Buildings,
city, marine, machines.}

\item[T6 4D replay] One clock drives geometry, appearance, camera and media; scrubbing
moves the whole scene. iTwin.js \texttt{RenderSchedule}; Samsara fuses GPS, multi-camera
video and sensors on one timeline where detected \emph{events}, not timestamps, are the
navigation anchors. \emph{All; critical for machines and mobility.}

\item[T7 Reality-capture fusion] Register as-built capture against as-designed model ---
point clouds where geometry must be measurable, Gaussian splats where fidelity is the
goal. \citet{mobilitydt2023} reconstruct vehicle geometry from crowdsourced imagery via
NeRF precisely because authoring is ``labour-intensive and time-consuming'', and report
PSNR falling 53--54\% outside ideal capture --- a caution against assuming capture
quality. \emph{Buildings, city, mobility.}

\item[T8 Video projection] Project a live feed onto the surfaces it sees, or lift
detections out of video into the scene as entities. Samsara Site Visibility is this
productised. \emph{Hospitals (privacy-constrained), machines, city.}

\item[T9 Ghost twin / residual] Render simulated and measured state in one frame --- a
translucent predicted pose against the measured one, or surfaces coloured by residual
rather than value. Essentially unserved commercially; Omniverse comes closest via USD
layers. Core to \citet{liao2025digital}. \textbf{Machines above all, and where DTaaS has
a genuine advantage.}

\item[T10 Linked abstraction levels] Schematic, plan, table and 3D as projections of one
selection and one clock. \citet{ali2024spatial} identify the general case as open
research. \emph{DTaaS shortcut: drive an embedded Grafana panel's \texttt{from}/%
\texttt{to} from the playhead.}

\item[T11 Streaming LOD] Hierarchical tiles, compression, per-element batching --- the
precondition for everything else at real scale. \emph{City and buildings critically.}

\item[T12 Attention routing] The scene navigates itself to the anomaly: camera frames
the element, the rest drops to x-ray. Almost entirely rules plus camera interpolation.
\emph{Hospitals and machines especially.}

\item[T13 Image-substrate overlay] A photograph replaces the model; readings, simulated
flows, queue lengths and predictions are drawn \emph{on the image} and animate with the
playhead. Three tiers, and choosing the cheapest sufficient one is most of the skill:
\textbf{(i) pixel anchoring} --- points and polygons in image coordinates, adequate
when the question is ``what is the value \emph{here}''; \textbf{(ii) ground-plane
homography} --- four image/world correspondences on a common plane yield a $3\times3$
homography, which is what makes simulation output directly projectable (SUMO positions,
crowd flow, predicted queue extent) in metres rather than pixels; \textbf{(iii) full
camera pose} --- intrinsics plus extrinsics for true 3D overlay, worth it only when
objects leave the ground plane. \emph{All domains; disproportionately valuable for city
and mobility.}

\item[T14 Live video beside simulation] Three arrangements: side-by-side on a shared
clock (video seeked to $t$ minus a declared offset); overlaid, which is \tech{13} with a
moving image; or a swipe divider between camera and matched-viewpoint render. The hard
part is \emph{time}, not display: stream latency runs from hundreds of milliseconds
(WebRTC) to seconds (HLS), so \req{11} demands the offset be declared, measurable and
adjustable. A side-by-side view with unstated skew is worse than no view, because
latency-induced divergence is indistinguishable from a wrong model. \emph{City,
mobility, machines, marine.}

\end{description}

\begin{keypoint}
\tech{13} should be built \emph{before} any 3D substrate. It exercises the whole binding
layer at a fraction of the cost, satisfies \req{9} by giving every twin some usable
view, and is what many users were asking for in the first place.
\end{keypoint}

\begin{longtable}{@{}L{0.8cm}L{4.0cm}C{1.6cm}C{1.6cm}C{1.6cm}C{1.6cm}C{1.6cm}@{}}
\caption{Technique applicability by domain.}\label{tab:domains}\\
\toprule
\textbf{ID} & \textbf{Technique} & \textbf{Hospitals} & \textbf{Buildings} &
\textbf{Machines} & \textbf{City/mob.} & \textbf{Marine} \\
\midrule
\endfirsthead
\toprule
\textbf{ID} & \textbf{Technique} & \textbf{Hospitals} & \textbf{Buildings} &
\textbf{Machines} & \textbf{City/mob.} & \textbf{Marine} \\
\midrule
\endhead
T1 & Semantic anchoring & Critical & Critical & Critical & Critical & Critical \\
T2 & Feature overrides & Critical & Critical & High & High & Medium \\
T3 & Surface field mapping & Critical & High & Low & Medium & High \\
T4 & Volumetric sim.\ fields & Medium & High & High & High & Critical \\
T5 & Flow advection & Medium & High & High & Critical & High \\
T6 & 4D replay & High & High & Critical & Critical & High \\
T7 & Reality-capture fusion & Medium & High & Medium & High & Medium \\
T8 & Video projection & High & Low & High & High & Medium \\
T9 & Ghost twin / residual & Medium & Medium & Critical & Medium & High \\
T10 & Linked abstraction levels & High & Medium & High & High & High \\
T11 & Streaming LOD & High & Critical & Low & Critical & Medium \\
T12 & Attention routing & Critical & High & High & High & Medium \\
T13 & Image-substrate overlay & High & High & High & Critical & Medium \\
T14 & Live video + simulation & High & Low & High & Critical & High \\
\bottomrule
\end{longtable}

\noindent
Three readings matter. \textbf{Hospitals and buildings are one rendering problem}: both
IFC-shaped, served by one adapter; what differs is data semantics, which lives in the
anchor and encoding layers. \textbf{Machines are genuinely different} --- no BIM library
helps with a pump or a robot cell --- but the scope is smaller because there is no
semantic model to honour, and \tech{9} pays off most there. \textbf{City, mobility and
marine share a renderer}, all georeferenced and leaning on \tech{4} and \tech{5}; both
\citet{rundel2023leveraging} and \citet{vasilijevic2024trondheim} independently arrive at
the same stack shape. And \tech{13} scores high in every column, because a photograph
exists long before a model does.


\section{Open-Source Packages}
\label{sec:packages}

Versions and licences verified against the npm registry and the GitHub API on 12--13
September 2026.

\begin{longtable}{@{}L{0.8cm}L{4.6cm}L{2.0cm}L{1.8cm}L{5.0cm}@{}}
\caption{Recommended package per technique (abridged).}\label{tab:packages}\\
\toprule
\textbf{ID} & \textbf{Package} & \textbf{Version} & \textbf{Licence} & \textbf{Role} \\
\midrule
\endfirsthead
\toprule
\textbf{ID} & \textbf{Package} & \textbf{Version} & \textbf{Licence} & \textbf{Role} \\
\midrule
\endhead
T1 & \pkg{@thatopen/components} & 3.4.8 & MIT & IFC property sets, classifier, lookup \\
T1 & \pkg{three-mesh-bvh} & 0.9.x & MIT & Fast picking; the inverse lookup \\
T2 & \pkg{three.js InstancedMesh} & --- & MIT & Per-instance colour, one draw call \\
T2 & \pkg{@thatopen/components-front} & 3.4.4 & MIT & Highlighter over fragments \\
T3/T4 & \pkg{vtk.js} & 32.x & BSD-3 & Interpolation, volumes, contours, streamlines \\
T3 & \pkg{deck.gl} & 9.4.0 & MIT & \texttt{HeatmapLayer}, plan-view equivalent \\
T5 & \pkg{deck.gl} \texttt{TripsLayer} & 9.4.0 & MIT & Time-windowed paths \\
T6 & \pkg{@reduxjs/toolkit} & 2.x & MIT & Playhead and control state (already present) \\
T7 & \pkg{Spark} & current & MIT & Gaussian-splat renderer for three.js \\
T7/T11 & \pkg{3DTilesRendererJS} & 0.4.x & Apache-2.0 & Streamed tiles and point clouds \\
T9 & --- & --- & --- & \textbf{No package exists}; built on T2 + T6 \\
T11 & \pkg{@thatopen/fragments} & 3.4.7 & MIT & Batched IFC, per-element addressing \\
T12 & \pkg{@react-three/drei} & 10.x & MIT & \texttt{CameraControls.fitToBox} \\
T13 & \pkg{react-konva} / \pkg{konva} & 19.2.7 / 10.5.0 & MIT & Glyph, polygon and label
layer over a photograph \\
T13 & \pkg{opencv-ts} (OpenCV.js) & 1.3.6 & MIT / Apache-2.0 & \texttt{findHomography},
reprojection error \\
T13 & \pkg{leaflet} + \pkg{react-leaflet} & 1.9.4 / 5.0.0 & BSD-2 / MIT &
\texttt{CRS.Simple} for very large plans \\
T14 & \pkg{hls.js} & 1.7.3 & Apache-2.0 & HLS where latency tolerance is seconds \\
T14 & \pkg{MediaMTX} & --- & MIT & \emph{Server-side}: RTSP\,$\rightarrow$\,WebRTC/HLS
gateway \\
\bottomrule
\end{longtable}

\paragraph{Renderer verdicts.} \textbf{Adopt} \pkg{react-three-fiber}+\pkg{drei} (MIT)
as \texttt{mesh}; \pkg{@thatopen/components} (MIT) as \texttt{aec} --- the only option
whose core abstraction is ``elements with queryable properties'', which is exactly a
telemetry binding, though it is on v3 and churns, so pin versions; \pkg{CesiumJS}
(Apache-2.0) as \texttt{geo}; \pkg{vtk.js} (BSD-3) as \texttt{field}, in its own canvas
with a synchronised camera. \textbf{Study but do not adopt} \pkg{iTwin.js} (MIT): the
entire technique catalogue already built, but a heavy platform-shaped dependency --- its
display-style / feature-override split is the cleanest model of the encode layer in
existence. \textbf{Reconsider later} \pkg{@pascal-app/*} (MIT) if in-client twin
\emph{authoring} is ever wanted. \textbf{Fallback} \pkg{@speckle/viewer} (Apache-2.0) if
building geometry arrives from many CAD tools rather than as IFC. \textbf{Reject}
\pkg{Babylon.js} (costs the R3F ecosystem, buys nothing), \pkg{GoJS} (2D and
commercial); \pkg{foxglove/studio} is archived and closed.

\paragraph{Transport and platform packages.} \pkg{mqtt} 5.x (MIT), \pkg{@stomp/stompjs}
7.x (Apache-2.0), \pkg{@influxdata/influxdb-client} 1.35.0 (MIT),
\pkg{@gitbeaker/rest} 43.8 (MIT, \emph{already shipped}), \pkg{react-iframe} 1.8.5 (MIT,
\emph{already shipped}), \pkg{comlink} 4.4.2 (Apache-2.0), and the ThingsBoard WebSocket
API, which needs no client library. Two of seven are already dependencies: the data
plane is mostly configuration, not new dependency surface.

\paragraph{Gaps and licence notes.} \tech{9} has no package and must be written --- about
a day of shader and encoder work, but only \emph{after} \tech{2} and \tech{6} exist.
\tech{8} has no maintained package (\pkg{three-projected-material} has not been pushed
since July 2024; write the $\approx$80-line projective shader instead). \tech{13} has no
integrated package: Konva draws, OpenCV solves, Leaflet pans, but nothing binds a live
signal to an image region along a playhead. And nothing at all covers Layers~1--4, which
is why they are the package's reason to exist. On licences: \pkg{xeokit-sdk} and
\pkg{realvirtual WEB} are AGPL-3.0 and excluded --- though realvirtual's GLB-plus-signal
wire format is worth studying; \pkg{Potree} and \pkg{loaders.gl} report
\texttt{NOASSERTION} and need checking; \pkg{konva} and \pkg{hls.js} report
\texttt{NOASSERTION} on GitHub while declaring MIT and Apache-2.0 on npm, so record the
exception explicitly rather than letting a scanner fail the build.


\section{Implementation Plan}
\label{sec:implementation}

The deliverable is a standalone npm package,
\pkg{@into-cps-association/dtaas-visualisation}, published beside the other DTaaS
packages. Standalone publication is deliberate: it forces the binding layer to be
independent of the client's routing, authentication and MUI theme, which is what makes
it reusable by other web-based DT platforms and testable without spinning up DTaaS.
Effort assumes one developer at $\approx$0.6~FTE; figures are planning estimates.

\begin{longtable}{@{}L{1.9cm}L{4.6cm}L{1.9cm}L{6.4cm}@{}}
\caption{Phases.}\label{tab:phases}\\
\toprule
\textbf{Phase} & \textbf{Deliverable} & \textbf{Techniques} & \textbf{Exit criterion} \\
\midrule
\endfirsthead
\toprule
\textbf{Phase} & \textbf{Deliverable} & \textbf{Techniques} & \textbf{Exit criterion} \\
\midrule
\endhead
\textbf{0} \newline 3 weeks & Layers 1--4 against the installed stack; \emph{no
substrate at all} & --- & Fixture replay emits the correct resolved encodings over a
scrubbed range, with no canvas on screen. If that is testable, \req{1} holds. \\
\addlinespace
\textbf{1} \newline 3 weeks & \texttt{image} adapter, calibration UI, Grafana embed &
T13, T10 & Live MQTT values drawn on a traffic or room photograph; scrub moves the
embedded Grafana panel. No 3D in the build. \\
\addlinespace
\textbf{2} \newline 4 weeks & \texttt{aec} adapter; ThingsBoard registry import & T1,
T2, T11 & A real building or hospital-floor twin recolouring by measured value through
the same binding layer. \\
\addlinespace
\textbf{3} \newline 2 weeks & Library asset type; binding editor committing via
gitbeaker & T12 & A second user reuses the asset by swapping substrate and anchors; a
threshold change lands as a merge request. \\
\addlinespace
\textbf{4} \newline 4 weeks & Replay over InfluxDB; ghost twin; \texttt{mesh} adapter &
T6, T9 & Scrubbing moves geometry, image overlay, table and residual together;
predicted and measured comparable in one frame. \\
\addlinespace
\textbf{5} \newline 3 weeks & \texttt{video} adapter; latency-offset control & T14, T8 &
Camera and simulation on one clock with skew displayed, adjustable and persisted. \\
\addlinespace
\textbf{6} \newline 4--6 weeks & \texttt{geo} and \texttt{field} adapters
(demand-driven) & T3, T4, T5, T7 & A georeferenced twin and a volumetric result render
through the Phase-0 binding layer. \\
\bottomrule
\end{longtable}

\begin{keypoint}
Phase~0 is the phase most likely to be cut under schedule pressure and the one that must
not be. Every later phase is cheap because of it, and neither the \texttt{channel} field
nor the playhead can be retrofitted without touching every module written after them.
\end{keypoint}

\paragraph{Integration points in the client.} A \texttt{visualise} sub-route under
\texttt{client/src/route/digitaltwins/}; \texttt{visualisation} registered as an asset
type in \texttt{route/library/} and the composition cart; a ``show on substrate''
affordance and shared playhead in \texttt{route/measurement/}; sibling VNC and web panes
in \texttt{route/workbench/}; a \texttt{visualisationSlice} in \texttt{client/src/store/}
holding playhead, connections, subscriptions, selection and layout --- with sample
buffers outside Redux; extended gitbeaker access in \texttt{client/src/util/}; and MIME
types plus range requests in \texttt{servers/lib}.

\paragraph{Principal risks.} Two are specific to the new substrates and both are failure
modes of \emph{credibility} rather than of code. \textbf{Homography error mistaken for
model error}: show reprojection error at calibration time, store it with the homography,
and refuse world-coordinate overlays above a threshold. \textbf{Video latency mistaken
for divergence} (high likelihood, high impact): declare the offset per feed, display it,
make it adjustable, and mark the pane ``unsynchronised'' when unknown. The others are
ordinary: \pkg{@thatopen/*} API churn (pin versions; confine contact to the
\texttt{aec} adapter); frame-rate collapse from per-sample Redux dispatch (buffers
outside Redux from Phase~0; CI performance test at 60~Hz and $10^4$ elements); bundle
size across seven adapters (lazy-load, CI size budget); InfluxDB query cost on wide
scrubs (window-and-cache, downsample, never query per frame); anchor drift on
re-export (persist \texttt{kind} and label, surface a repair UI on failure); licence
contamination (CI allowlist); and privacy in hospital and public-space video (restricted
channels, default off, documented lawful basis --- \citet{ali2024spatial} note location
data alone can identify an individual).

\paragraph{What to build first.} If one sentence survives contact with the schedule:
\textbf{build the transport, the hot/cold time store and the \texttt{SubstrateAdapter}
interface against the RabbitMQ and InfluxDB already deployed, then prove them on a
photograph} --- not on a 3D model. That route reaches a user-recognisable result in
about six weeks, needs no WebGL, exercises every layer the 3D substrates will later
reuse, and leaves a fallback view for every twin that never acquires a model.


\section{Evaluation}
\label{sec:evaluation}

\citet{mobilitydt2023} separate \emph{fidelity} (precision independent of latency) from
\emph{physical-to-digital synchronicity} (discrepancy caused by latency); the separation
is the useful move, because the two trade against each other and one ``quality'' number
hides the trade. Their own measurement makes the point: compressing offloaded imagery
from $>18$~MB to $\approx6$~MB cost less than 0.1\% of PSNR while cutting latency
substantially. The analogues here are \emph{encoding fidelity} (rendered property versus
the encoding evaluated in the pure layer --- an automatable readback test) and
\emph{synchronicity} (p50/p95/p99 publish-to-pixel latency at a stated element count).

Three further measures follow from the new substrates and from \tech{9}.
\textbf{Residual credibility}: report the residual's own error budget beside it, since a
ghost twin showing 2\,\textdegree C divergence is useless if pipeline uncertainty is
3\,\textdegree C --- \citet{liao2025digital} model this discipline with quantitative
metrics rather than visual plausibility. \textbf{Registration accuracy}: homography
reprojection error, expressed in ground-plane metres as well as pixels, because three
pixels near the horizon of a traffic photograph can be several metres.
\textbf{Synchronisation skew}: the signed video-frame-to-playhead difference, reported
continuously since transport latency drifts; the question is not whether skew is zero
but whether it is known, bounded and smaller than the phenomenon being compared. Finally,
\citet{chen2024tangible} evaluate an MR airport digital tower with ten licensed
controllers on workload, situational awareness and trust --- the right constructs for an
operator-facing view, and not substitutable by usability, since a view can be easy to use
and still degrade situational awareness by drawing attention to the wrong thing. All but
the operator study belong in CI from the phase that introduces them.

\section{Conclusion}
\label{sec:conclusion}

\subsection{Findings}

\begin{enumerate}[itemsep=2pt]
  \item The DTaaS client has no spatial visualisation capability today beyond embedded
        time series.
  \item Every surveyed commercial platform shares one architecture --- anchor, encode,
        sustain --- and sells a no-code editor over the encode layer.
  \item The renderer is therefore the wrong thing to standardise on. Transport
        normalisation, the temporal store, anchor resolution and the encoding vocabulary
        are the durable artefacts, and none exist off the shelf.
  \item Not every substrate is a 3D scene. Photographs and live video are first-class
        places to put simulation output; generalising the adapter interface costs
        nothing and changes what ships in the first two months.
  \item Roughly a third of the layer is already deployed --- historian, broker, device
        registry, charting, versioned storage and identity need to be \emph{used}, not
        built.
  \item Hospitals and buildings are one 3D problem, machines another, city/mobility/marine
        a third: three geometric adapters plus \texttt{image}, \texttt{video} and
        \texttt{embed}.
  \item Simulation--measurement co-visualisation is open ground, and DTaaS is unusually
        placed to take it.
  \item One early decision governs half the catalogue: every encoding must read
        \texttt{valueAt(playhead)}, never ``latest value''.
  \item Licence discipline removes two of the strongest candidates
        (\pkg{xeokit-sdk}, \pkg{realvirtual WEB}) and flags two more.
\end{enumerate}

\subsection{Positioning and next step}

The literature converges on this architecture from several directions.
\citet{rundel2023leveraging} show that the escape from domain-rebuild cost is to derive
the scene from data that already exists --- and a photograph is data that already exists.
\citet{vasilijevic2024trondheim} adopt a deliberately two-tier visualisation, Grafana for
series and a game-engine/Cesium layer for geospatial rendering, which is almost exactly
the split proposed here between the \texttt{embed} adapter and the geometric substrates,
coupled through a shared clock rather than merged. \citet{ali2024spatial} caution that
\tech{10} is not merely engineering. \citet{mobilitydt2023} supply the evaluation
constructs and, by measuring how expensive geometry reconstruction is, the strongest
argument for using imagery where it suffices. \citet{liao2025digital} demonstrate the
ghost twin in miniature; \citet{chen2024tangible} supply the human-factors constructs and
the case for shared, multi-operator views; and \citet{yang2024digital} identify data
integration and quality, not rendering, as the binding constraint on DT adoption ---
consistent with the claim that the binding layer is the product.

\medskip
\noindent\textbf{Next step.} Build Phase~0 against the services already running, then
Phase~1 on a photograph. Six weeks to something a user recognises as their own twin, no
WebGL required, and the whole binding layer proven before a single 3D asset is
commissioned.


\appendix
\section{Source Provenance and Licence Register}
\label{app:sources}

\subsection{Supplied sources}

Eight of the eleven files in \texttt{sources/} are papers and were text-extracted
successfully: \citet{yang2024digital} (\emph{Digital Twins in Construction}, Buildings
14(9):2616); \citet{ali2024spatial} (\emph{Spatial Digital Twins}, PFG 92:761--778);
\citet{vasilijevic2024trondheim} (\emph{A Digital Twin of the Trondheim Fjord}, JMSE
12(9):1530); \citet{chen2024tangible} (\emph{Tangible digital twin with shared
visualization for collaborative air traffic management operations}, Transportation
Research Part~C); \citet{liao2025digital} (\emph{An augmented reality-enabled digital
twin system for reconfigurable soft robots}, Computers in Industry);
\citet{rundel2023leveraging} (\emph{Leveraging digital twin and game-engine for traffic
simulations and visualizations}, Frontiers in Virtual Reality 4:1048753); and
\citet{mobilitydt2023} (\emph{Visualization of Mobility Digital Twin}, IEEE).

Three caveats are recorded rather than hidden. \fn{1-s2.0-S0278612525000834-main.pdf}
uses CID-encoded embedded fonts and could not be text-extracted by either the direct
extractor or Word conversion; it is identified from metadata as \citet{iiot2025llm} and
\emph{nothing in this report relies on its content}. And \fn{dam-visualisation.pdf},
\fn{main.pdf} and \fn{main-2.pdf} (each $\approx$61~KB) are \textbf{not papers} --- each
is a saved ScienceDirect ``JavaScript is disabled / Request Verification'' error page.

\begin{keypoint}
If \fn{dam-visualisation.pdf} was meant to contribute a dam or hydraulic-structure case
study --- a domain not otherwise covered --- it should be re-downloaded and this report
extended, since that domain would exercise \tech{4} and \tech{5} heavily.
\end{keypoint}

\subsection{Extraction and verification method}

No LaTeX distribution, Python, Node.js or Poppler was available when the sources were
read. PDF text was recovered with a purpose-written PowerShell routine that walks each
file's stream objects, inflates \texttt{FlateDecode} payloads via
\texttt{System.IO.Compression.DeflateStream}, and extracts text-showing operators
(\texttt{Tj}, \texttt{TJ}) with PDF string-escape decoding. Inter-token spacing in the
recovered text is imperfect but word content is faithful; quotations were re-spaced by
hand and checked against context. Package versions and licences were read from the npm
registry and the GitHub REST API on 12--13 September 2026; where the API reports
\texttt{NOASSERTION} --- no determinable SPDX identifier --- that is stated rather than
guessed.

\subsection{Licence register}

\begin{longtable}{@{}L{5.6cm}L{3.0cm}L{6.0cm}@{}}
\toprule
\textbf{Package} & \textbf{Licence} & \textbf{Status for DTaaS} \\
\midrule
\endhead
\pkg{react-three-fiber}, \pkg{drei}, \pkg{three.js}, \pkg{three-mesh-bvh} & MIT & Clear \\
\pkg{@thatopen/components}, \pkg{-front}, \pkg{fragments} & MIT & Clear \\
\pkg{deck.gl}, \pkg{Spark}, \pkg{SuperSplat} & MIT & Clear \\
\pkg{iTwin.js}, \pkg{@pascal-app/*} & MIT & Clear; not adopted (see
\S\ref{sec:packages}) \\
\pkg{react-konva}, \pkg{react-leaflet}, \pkg{opencv-ts} & MIT & Clear \\
\pkg{mqtt}, \pkg{@influxdata/influxdb-client} & MIT & Clear \\
\pkg{@gitbeaker/rest}, \pkg{react-iframe} & MIT & Clear; already client dependencies \\
\pkg{MediaMTX} & MIT & Clear; deployment component, not a client dependency \\
\pkg{CesiumJS}, \pkg{3DTilesRendererJS} & Apache-2.0 & Clear \\
\pkg{@speckle/viewer}, \pkg{@stomp/stompjs}, \pkg{comlink} & Apache-2.0 & Clear \\
OpenCV (behind \pkg{opencv-ts}) & Apache-2.0 & Clear \\
\pkg{vtk.js} & BSD-3-Clause & Clear \\
\pkg{leaflet} & BSD-2-Clause & Clear \\
\pkg{konva} & MIT (npm); \texttt{NOASSERTION} (GitHub) & Clear; record as a CI
exception \\
\pkg{hls.js} & Apache-2.0 (npm); \texttt{NOASSERTION} (GitHub) & Clear; record as a CI
exception \\
\pkg{Potree}, \pkg{loaders.gl} & \texttt{NOASSERTION} & \textbf{Check before use} \\
\pkg{xeokit-sdk}, \pkg{realvirtual WEB} & \textbf{AGPL-3.0} & \textbf{Excluded} ---
network copyleft \\
\pkg{GoJS} & Commercial & \textbf{Excluded} \\
\bottomrule
\end{longtable}


\bibliographystyle{plainnat}
\bibliography{references}

\end{document}
